\section{Common analytic functions}\label{sec:analysis}

For each auxiliary index \(a\) and column polynomial \(P\), define
\[
 f_{a,P}(z,z')=[u^a]P(e^z,e^{z'},(z'-\theta z+u_i)_{i=1}^m).
\]
For a monomial column this is
\(\binom\alpha a e^{hz+kz'}(z'-\theta z)^{|\alpha|-|a|}\), or
zero if \(a\not\leq\alpha\).
At center \(j\) set
\(z=j\log2+\log(1+t)\), \(z'=j\log3+\log(1+s)\).
Writing \(\delta_i=\theta-r_i\), let
\[
 B_i=G_i(s)-\log(1+s)-r_i(G_i(t)-\log(1+t)),\quad
 h_i=\delta_i\log(1+t)+B_i.
\]
The exact identity \(G_i(s)-r_iG_i(t)=z'-\theta z+h_i\)
has no center-dependent error.

\begin{lemma}\label{lem:translation}
Every original row is a finite sum of scalar multiples of the common
coefficient tests
\([t^{\ell'_1}s^{\ell'_2}]f_{a,P}(z,z')\), with
\(a\geq\beta\), \(wa\leq H\), and \(\ell'_1+\ell'_2<H/v_0\).
The scalars are independent of \(P\). Their number per row is at most
\[
 C_H=(\lfloor H\rfloor+1)^m
                  (\lfloor H/v_0\rfloor+1)^2,
\]
and each is bounded by \((\lfloor H\rfloor+1)^m\) times
\begin{equation}\label{eq:translation}
 \exp\{-\nu w(a-\beta)+HE_{\rm tr}\},\qquad
 E_{\rm tr}=\frac\nu F+\frac{\log2}{v_0}
                       +\frac{\nu+\log12}{w_*}.
\end{equation}
\end{lemma}
\begin{proof}
The scalar in the convolution with base indices \(d,e\) is
\[
 \binom a\beta[t^ds^e]\prod_i h_i^{a_i-\beta_i},
 \qquad 0\leq d\leq\ell_1,\quad0\leq e\leq\ell_2.
\]
In its expansion let \(b_i\) factors be tails \(B_i\). A nonzero
contribution requires \(\sum_iT_ib_i\leq d+e<H/v_0\), hence
\(wb<H/F\). Because \(w_i-1<\log q_i\leq w_i\), its approximation
factors are bounded by
\(\exp(-\nu w(a-\beta)+\nu H/F+\nu H/w_*)\).
On \(|t|=|s|=1/2\), logarithms have modulus at most \(\log2<1\),
and \(|B_i|\leq3\). The two binomial factors cost at most
\(4^{H/w_*}\), the tails at most \(3^{H/w_*}\), and Cauchy's
coefficient estimate at most \(2^{H/v_0}\). This proves the bound
after summing at most \((\lfloor H\rfloor+1)^m\) tail choices.
Indeed, each auxiliary coordinate of \(a\) is at most \(\lfloor H\rfloor\),
since \(w_i\geq1\). For a fixed \(a\), the convolution has exactly
\((\ell_1+1)(\ell_2+1)\) possible splits, which is at most
\((\lfloor H/v_0\rfloor+1)^2\). Thus the combined counting and
scalar overhead per row is bounded by
\[
 Q_H=(\lfloor H\rfloor+1)^{2m}
                  (\lfloor H/v_0\rfloor+1)^2.
\]
All data other than \(H\) are fixed, so \(\log Q_H=O(\log H)=o(H)\).
\end{proof}
